% ══════════════════════════════════════════════════════════════════════════════════════════
%  Diccionario humano de citoquinas: práctica de Inmunología (diapositivas, v3, español)
%  Compilar:  python slides/v3/make_figures.py es   (figuras, números y tablas desde el notebook)
%             cd slides/v3_es && latexmk -xelatex main.tex
% ══════════════════════════════════════════════════════════════════════════════════════════
\documentclass[aspectratio=169,10pt]{beamer}
\usetheme[progressbar=frametitle,numbering=fraction,sectionpage=progressbar]{metropolis}

\usepackage{iftex}
\ifPDFTeX\usepackage[T1]{fontenc}\usepackage[utf8]{inputenc}\fi
\usepackage{polyglossia}
\setmainlanguage{spanish}
\usepackage{amsmath,amssymb}
\usepackage{booktabs}
\usepackage{array}
\usepackage{graphicx}
\usepackage{xcolor}
\usepackage{listings}
\usepackage{url}
\usepackage{tikz}
\usetikzlibrary{arrows.meta,positioning,calc,fit,backgrounds,shapes.geometric,decorations.pathreplacing,matrix}
\usepackage{appendixnumberbeamer}
\graphicspath{{img/}}

\ifXeTeX
  \setsansfont[Extension=.otf,UprightFont=*-Regular,BoldFont=*-Bold,ItalicFont=*-Italic,BoldItalicFont=*-BoldItalic]{FiraSans}
  \setmonofont[Extension=.otf,UprightFont=*-Regular,BoldFont=*-Bold,Scale=0.92]{FiraMono}
\fi
\newcommand{\numnsamples}{4\,224}
\newcommand{\numngenesraw}{18\,120}
\newcommand{\numngenes}{15\,535}
\newcommand{\numngenesremoved}{2\,585}
\newcommand{\numumiMono}{6\,747}
\newcommand{\numumiNK}{3\,147}
\newcommand{\numpbsratio}{5.3}
\newcommand{\numtoptwentycpm}{35\%}
\newcommand{\numtoptwentylog}{1.6\%}
\newcommand{\numpcaone}{66}
\newcommand{\numpcatwo}{17}
\newcommand{\numpcathree}{3.4}
\newcommand{\numpcafirstfour}{88}
\newcommand{\numntested}{12\,743}
\newcommand{\numnup}{1\,353}
\newcommand{\numndown}{1\,282}
\newcommand{\numnrawsig}{7\,239}
\newcommand{\numnadjsig}{6\,585}
\newcommand{\numnexpectedfp}{637}
\newcommand{\numnnullraw}{750}
\newcommand{\numnnulladj}{0}
\newcommand{\numtobs}{10.6}
\newcommand{\numpobs}{3.9\times10^{-7}}
\newcommand{\numlfcobs}{3.66}
\newcommand{\numfoldobs}{12.6}
\newcommand{\numkstar}{6\,585}
\newcommand{\nummediancpm}{$-0.09$}
\newcommand{\nummediandeseq}{$+0.00$}
\newcommand{\numndegdeseq}{2\,635}
\newcommand{\numndegt}{2\,198}
\newcommand{\numndegboth}{1\,997}
\newcommand{\numreffect}{0.97}
                                    % numbers computed by make_figures.py

% ── colours (same hues as the notebook figures) ─────────────────────────────────────────
\definecolor{ucoblue}{HTML}{184F95}
\definecolor{bluemid}{HTML}{2A78D6}
\definecolor{accentred}{HTML}{E34948}
\definecolor{greyink}{HTML}{52514E}
\definecolor{softgrey}{HTML}{F0EFEC}
\definecolor{orangeslot}{HTML}{EB6834}
\definecolor{greenslot}{HTML}{1BAF7A}
\definecolor{yellowslot}{HTML}{EDA100}
\definecolor{magentaslot}{HTML}{E87BA4}
\setbeamercolor{frametitle}{bg=ucoblue,fg=white}
\setbeamercolor{progress bar}{fg=accentred,bg=ucoblue!30}
\setbeamercolor{alerted text}{fg=accentred}
\setbeamercolor{block title}{bg=ucoblue!85,fg=white}
\setbeamercolor{block body}{bg=softgrey,fg=black}
\setbeamercolor{title separator}{fg=accentred}
\setbeamercolor{section title}{fg=ucoblue}
\setbeamerfont{block title}{size=\small}
\setbeamertemplate{navigation symbols}{}

% ── helpers ───────────────────────────────────────────────────────────────────────────
\newcommand{\gene}[1]{\textit{#1}}
\newcommand{\code}[1]{\texttt{#1}}
\newcommand{\IFNg}{IFN-$\gamma$}
\newcommand{\figfit}[2][0.66]{\includegraphics[width=\linewidth,height=#1\textheight,keepaspectratio]{#2}}
\newcommand{\lfc}{\log_2\mathrm{FC}}
\newcommand{\timebox}[1]{\hfill{\footnotesize\textcolor{greyink}{$\circ$ #1}}}

\lstset{basicstyle=\ttfamily\scriptsize,keywordstyle=\color{bluemid}\bfseries,commentstyle=\color{greyink}\itshape,
        stringstyle=\color{accentred},showstringspaces=false,frame=single,rulecolor=\color{ucoblue!40},
        backgroundcolor=\color{softgrey},language=Python,columns=fullflexible,keepspaces=true,
        literate={→}{{$\rightarrow$}}1}

\title{Diccionario humano de citoquinas}
\subtitle{De los conteos crudos a la expresión diferencial\\ AnnData $\cdot$ normalización $\cdot$ PCA $\cdot$ análisis diferencial}
\author{Martin Garrido-Rodriguez}
\institute{UCO $\cdot$ Práctica de Inmunología}
\date{Basado en Oesinghaus \emph{et al.} (2025), \emph{A single-cell cytokine dictionary of human peripheral blood}}

\begin{document}

\maketitle

% ════════════════════════════════════════════════════════════════════════════════════════
\begin{frame}{Plan de la práctica: 2 h 30 min}
  \centering
  \begin{tikzpicture}[font=\scriptsize, every node/.style={align=center},
      seg/.style={draw=white,thick,minimum height=1.3cm,text=white,font=\scriptsize\bfseries,inner sep=1pt}]
    \def\u{0.0866}   % cm per minute, total width = 150 min = 13 cm
    \node[seg,fill=ucoblue,       minimum width={25*\u cm}, anchor=west] at (0,0)   {1\\ Intro};
    \node[seg,fill=bluemid,       minimum width={40*\u cm}, anchor=west] at (25*\u,0)  {2\\ Datos y normalización};
    \node[seg,fill=greyink!60,    minimum width={10*\u cm}, anchor=west] at (65*\u,0)  {\tiny pausa};
    \node[seg,fill=orangeslot,    minimum width={35*\u cm}, anchor=west] at (75*\u,0)  {3\\ PCA};
    \node[seg,fill=accentred,     minimum width={35*\u cm}, anchor=west] at (110*\u,0) {4\\ Expr. diferencial};
    \node[seg,fill=greyink,       minimum width={5*\u cm},  anchor=west] at (145*\u,0) {};
    \foreach \t/\x in {0/0,25/25,65/65,75/75,110/110,150/150}
      \node[below=0.12cm,text=greyink] at (\x*\u,-0.65) {\t'};
  \end{tikzpicture}

  \vspace{0.45cm}
  \small
  \begin{tabular}{@{}p{0.07\linewidth}p{0.45\linewidth}p{0.25\linewidth}p{0.12\linewidth}@{}}
    \toprule
    \textbf{Parte} & \textbf{Tema} & \textbf{Formato} & \textbf{Tiempo}\\\midrule
    \textbf{1} & Introducción: Python, Jupyter, Colab, RNA-seq, célula única, pseudobulk, AnnData & diapositivas & 25 min\\[2pt]
    \textbf{2} & Manipular y normalizar los datos & notebook & 40 min\\[2pt]
    & \textit{Descanso} & & 10 min\\[2pt]
    \textbf{3} & Análisis de Componentes Principales & diapositivas + notebook & 35 min\\[2pt]
    \textbf{4} & Expresión diferencial: citoquina vs control & diapositivas + notebook & 35 min\\[2pt]
    & Cierre & & 5 min\\
    \bottomrule
  \end{tabular}
\end{frame}

% ════════════════════════════════════════════════════════════════════════════════════════
\section{Parte 1: Introducción}

\begin{frame}[fragile]{¿Qué es Python?}
  \begin{columns}[T]
    \begin{column}{0.52\linewidth}
      \begin{itemize}\setlength\itemsep{4pt}
        \item Un \alert{lenguaje de programación de propósito general}: gratuito, legible y con una enorme comunidad.
        \item Describes \emph{qué hacer} paso a paso y el ordenador lo hace, con 10 muestras o con 10 millones.
        \item En ciencia es el \alert{pegamento} entre los datos y el análisis. Paquetes que usamos hoy:
      \end{itemize}
      \small
      \begin{tabular}{@{}ll@{}}
        \code{numpy} & matrices y álgebra lineal\\
        \code{pandas} & tablas\\
        \code{scipy}, \code{statsmodels} & estadística\\
        \code{matplotlib} & gráficos\\
        \code{scikit-learn} & aprendizaje automático\\
        \code{anndata} & contenedores de datos ómicos\\
      \end{tabular}
    \end{column}
    \begin{column}{0.46\linewidth}
\begin{lstlisting}
genes  = ["STAT1", "CXCL9", "ACTB"]
counts = {"STAT1": 73853, "ACTB": 36500}
total  = sum(counts.values())
print(f"{total:,} UMI en total")
\end{lstlisting}
      \vspace{0.15cm}
      \begin{block}{Hoy no necesitas ser programador}
        Te damos el código. Tu trabajo: \textbf{ejecutarlo, leer el resultado y entender qué significan los números}.
      \end{block}
    \end{column}
  \end{columns}
\end{frame}

\begin{frame}{¿Qué es un notebook de Jupyter?}
  \begin{columns}[c]
    \begin{column}{0.56\linewidth}
      \centering
      \begin{tikzpicture}[font=\scriptsize, every node/.style={align=left}]
        \node[draw=greyink!60,fill=white,minimum width=6.6cm,minimum height=3.9cm,anchor=north west] (page) at (0,0) {};
        % markdown cell
        \node[draw=none,fill=softgrey,minimum width=6.2cm,minimum height=0.95cm,anchor=north west] (md) at (0.2,-0.2) {};
        \node[anchor=north west,font=\small\bfseries] at (0.35,-0.28) {Parte 2: normalización};
        \node[anchor=north west,font=\scriptsize,text=greyink] at (0.35,-0.68) {Los conteos crudos no son comparables\ldots};
        % code cell
        \node[draw=bluemid,fill=bluemid!6,minimum width=6.2cm,minimum height=0.95cm,anchor=north west] (code) at (0.2,-1.3) {};
        \node[anchor=north west,font=\ttfamily\scriptsize] at (0.35,-1.42) {adata = ad.read\_h5ad("data.h5ad")};
        \node[anchor=north west,font=\ttfamily\scriptsize] at (0.35,-1.78) {print(adata.shape)};
        \node[circle,fill=bluemid,text=white,inner sep=2pt,font=\tiny\bfseries] at (0.02,-1.8) {$\blacktriangleright$};
        % output
        \node[anchor=north west,font=\ttfamily\scriptsize,text=ucoblue] at (0.35,-2.45) {(4224, 18120)};
        \node[anchor=north west,font=\scriptsize,text=greyink] at (3.3,-2.45) {$\leftarrow$ el resultado};
        % labels
        \node[anchor=west,text=accentred,font=\scriptsize\bfseries] at (6.9,-0.65) {texto};
        \node[anchor=west,text=accentred,font=\scriptsize\bfseries] at (6.9,-1.75) {código};
        \node[anchor=west,text=accentred,font=\scriptsize\bfseries] at (6.9,-2.55) {resultado};
        \draw[accentred,-{Stealth}] (6.85,-0.65) -- (6.45,-0.65);
        \draw[accentred,-{Stealth}] (6.85,-1.75) -- (6.45,-1.75);
        \draw[accentred,-{Stealth}] (6.85,-2.55) -- (6.45,-2.55);
      \end{tikzpicture}
    \end{column}
    \begin{column}{0.42\linewidth}\small
      \begin{itemize}\setlength\itemsep{3pt}
        \item Un documento que mezcla \alert{texto, código y resultados}.
        \item Dos tipos de celdas: de \textbf{texto} (Markdown) y de \textbf{código} (Python).
        \item Ejecutar una celda: \alert{Shift\,+\,Enter}. El resultado aparece debajo.
        \item Un \textbf{kernel} (el proceso de Python que hay detrás) \alert{recuerda} las variables: ejecuta las celdas \textbf{en orden}.
        \item Reiniciar el kernel $\Rightarrow$ se borra la memoria.
      \end{itemize}
    \end{column}
  \end{columns}
\end{frame}

\begin{frame}{¿Cómo lo vamos a ejecutar? Google Colab}
  \begin{columns}[c]
    \begin{column}{0.60\linewidth}
      \centering
      \begin{tikzpicture}[font=\scriptsize, every node/.style={align=left}]
        \draw[rounded corners=3pt,draw=greyink!70,fill=white] (0,0) rectangle (7.6,-4.2);
        \fill[greyink!12] (0,0) rectangle (7.6,-0.55);
        \node[anchor=west,font=\scriptsize\bfseries] at (0.15,-0.28) {Practica\_v3\_es.ipynb};
        \node[anchor=east,draw=bluemid,fill=bluemid!10,rounded corners=2pt,inner sep=2pt,font=\tiny] at (7.45,-0.28) {Conectar};
        \fill[greyink!5] (0,-0.55) rectangle (7.6,-0.95);
        \node[anchor=west,font=\tiny] at (0.15,-0.75) {Archivo\ \ Editar\ \ Ver\ \ Insertar\ \ \textbf{Entorno}\ \ Herram.\ \ Ayuda};
        \node[anchor=west,font=\tiny,draw=greyink!50,inner sep=1.5pt] at (4.7,-0.75) {+ Código};
        \node[anchor=west,font=\tiny,draw=greyink!50,inner sep=1.5pt] at (5.75,-0.75) {+ Texto};
        \fill[greyink!8] (0,-0.95) rectangle (1.45,-4.2);
        \node[anchor=north west,font=\tiny] at (0.1,-1.05) {$\equiv$ Índice\\ Parte 1\\ Parte 2\\ \quad 2.1 \ldots\\ Parte 3\\ Parte 4};
        \node[draw=bluemid,fill=bluemid!6,minimum width=5.6cm,minimum height=1.15cm,anchor=north west] at (1.75,-1.15) {};
        \node[anchor=north west,font=\ttfamily\tiny] at (2.2,-1.25) {adata = ad.read\_h5ad(DATA\_PATH)};
        \node[circle,fill=bluemid,text=white,inner sep=2.5pt,font=\tiny] at (1.95,-1.7) {$\blacktriangleright$};
        \node[anchor=north west,font=\ttfamily\tiny,text=ucoblue] at (2.2,-1.75) {\checkmark\ Cargado en memoria};
        \node[draw=none,fill=softgrey,minimum width=5.6cm,minimum height=0.8cm,anchor=north west] at (1.75,-2.6) {};
        \node[anchor=north west,font=\scriptsize\bfseries] at (1.9,-2.68) {2.2 Un número cada vez};
        \node[anchor=north west,font=\tiny,text=greyink] at (1.9,-3.02) {las celdas de texto explican, las de código calculan};
        % numbered callouts
        \foreach \n/\x/\y/\tx/\ty in {1/7.2/0.28/7.2/-0.12, 2/1.55/-1.0/1.55/-1.3, 3/3.9/0.28/3.9/-0.55, 4/-0.3/-1.6/0.05/-1.6} {
          \node[circle,fill=accentred,text=white,inner sep=1.5pt,font=\tiny\bfseries] (c\n) at (\x,\y) {\n};
          \draw[accentred,thick] (c\n) -- (\tx,\ty);
        }
      \end{tikzpicture}
    \end{column}
    \begin{column}{0.38\linewidth}\small
      \begin{itemize}\setlength\itemsep{3pt}
        \item Jupyter en tu \alert{navegador}: nada que instalar. Gratis, con una cuenta de Google.
        \item Tu código se ejecuta en un \alert{ordenador temporal} en la nube de Google.
      \end{itemize}
      \vspace{0.1cm}
      \begin{enumerate}\setlength\itemsep{2pt}\footnotesize
        \item[\textcolor{accentred}{\textbf{1}}] \textbf{Conectar} a una sesión
        \item[\textcolor{accentred}{\textbf{2}}] $\blacktriangleright$ \textbf{ejecutar} una celda
        \item[\textcolor{accentred}{\textbf{3}}] menú \textbf{Entorno de ejecución}: ejecutar todo, reiniciar
        \item[\textcolor{accentred}{\textbf{4}}] \textbf{Índice}: saltar entre partes
      \end{enumerate}
    \end{column}
  \end{columns}
\end{frame}

\begin{frame}{Kit de supervivencia de Colab}
  \begin{columns}[T]
    \begin{column}{0.52\linewidth}
      \small
      \begin{tabular}{@{}>{\ttfamily}p{0.38\linewidth}p{0.58\linewidth}@{}}
        \toprule
        \textrm{\textbf{Atajo}} & \textbf{Qué hace}\\\midrule
        Shift+Enter & ejecuta y pasa a la siguiente\\
        Ctrl+Enter & ejecuta y se queda\\
        Ctrl+M B & inserta una celda debajo\\
        Ctrl+M D & borra la celda\\
        \bottomrule
      \end{tabular}
      \vspace{0.3cm}

      \textbf{Si algo se rompe}
      \begin{itemize}\footnotesize\setlength\itemsep{2pt}
        \item \textbf{Entorno de ejecución $\rightarrow$ Reiniciar sesión} y vuelve a ejecutar desde el principio.
        \item ¿Te has quedado atrás en la Parte 2? Usa la \alert{celda de rescate} al inicio de la Parte 3.
      \end{itemize}
    \end{column}
    \begin{column}{0.44\linewidth}
      \begin{block}{Conviene saber}
        \footnotesize
        \begin{itemize}\setlength\itemsep{3pt}
          \item Las sesiones son \alert{temporales}: archivos y variables desaparecen al terminar. Descargamos los datos al principio (alrededor de un minuto).
          \item Una sesión inactiva se desconecta: reconecta y vuelve a ejecutar.
          \item Las celdas cuyo título termina en \emph{``ejecutar una vez, no hace falta leer''} están plegadas a propósito: son código auxiliar para los gráficos.
          \item Para conservar tus cambios: \textbf{Archivo $\rightarrow$ Guardar una copia en Drive}.
        \end{itemize}
      \end{block}
    \end{column}
  \end{columns}
\end{frame}

% ────────────────────────────────────────────────────────────────────────────────────────
\begin{frame}{¿Qué es la transcriptómica mediante RNA-seq?}
  \centering
  \begin{tikzpicture}[font=\scriptsize, every node/.style={align=center},
      stp/.style={draw=ucoblue,fill=ucoblue!8,rounded corners,minimum height=1.35cm,minimum width=2.05cm,inner sep=2pt}]
    \node[stp] (a) {tejido o\\ células};
    \node[stp,right=0.35cm of a] (b) {\textbf{extracción}\\ de RNA};
    \node[stp,right=0.35cm of b] (c) {\textbf{librería}\\ fragmentos +\\ adaptadores};
    \node[stp,right=0.35cm of c] (d) {\textbf{secuenciación}\\ millones de\\ lecturas};
    \node[stp,right=0.35cm of d] (e) {\textbf{alineamiento}\\ lecturas $\to$ genes};
    \node[stp,right=0.35cm of e,fill=accentred!10,draw=accentred] (f) {\textbf{matriz de}\\ \textbf{conteos}};
    \foreach \x/\y in {a/b,b/c,c/d,d/e,e/f} \draw[-{Stealth},thick] (\x) -- (\y);
  \end{tikzpicture}
  \vspace{0.35cm}

  \begin{columns}[T]
    \begin{column}{0.5\linewidth}\small
      \begin{itemize}\setlength\itemsep{3pt}
        \item El \alert{transcriptoma} es el conjunto de todas las moléculas de RNA de una muestra: qué genes están \emph{encendidos} y cuánto.
        \item El \textbf{RNA-seq bulk} mide \alert{un perfil promedio por muestra}: todas las células están mezcladas.
        \item Un \textbf{conteo} es el número de moléculas (o lecturas) asignadas a un gen.
      \end{itemize}
    \end{column}
    \begin{column}{0.46\linewidth}
      \centering\footnotesize
      \begin{tabular}{lrrr}
        \toprule
        gen & muestra 1 & muestra 2 & muestra 3\\\midrule
        \gene{STAT1} & 5\,310 & 12\,874 & 4\,990\\
        \gene{CXCL9} & 12 & 1\,950 & 9\\
        \gene{ACTB} & 36\,500 & 41\,200 & 33\,800\\
        $\vdots$ & $\vdots$ & $\vdots$ & $\vdots$\\
        \bottomrule
      \end{tabular}\\[2pt]
      \scriptsize una matriz de conteos: \alert{genes $\times$ muestras} (números ilustrativos)
    \end{column}
  \end{columns}
\end{frame}

\begin{frame}{RNA-seq bulk y de célula única: los dos flujos de trabajo}
  \begin{columns}[c]
    \begin{column}{0.3\linewidth}
      \centering\includegraphics[width=\linewidth,height=0.8\textheight,keepaspectratio]{tzec_fig1_bulk.png}
    \end{column}
    \begin{column}{0.3\linewidth}
      \centering\includegraphics[width=\linewidth,height=0.8\textheight,keepaspectratio]{tzec_fig2_sc.png}
    \end{column}
    \begin{column}{0.36\linewidth}\small
      \begin{itemize}\setlength\itemsep{4pt}
        \item La misma lógica con distinta unidad: una \alert{muestra} (bulk) o una \alert{célula} (célula única).
        \item Ambos terminan en una \alert{matriz de conteos}, que luego se normaliza y se analiza (expresión diferencial, enriquecimiento\ldots).
        \item La célula única añade \alert{códigos de barras} para asignar cada lectura a su célula, y control de calidad por célula.
      \end{itemize}
      {\tiny Figuras 1 y 2 de Tzec-Interi\'an \emph{et al.}\ 2025, \emph{Quant.\ Biol.}\ (CC BY)}
    \end{column}
  \end{columns}
\end{frame}

\begin{frame}{Transcriptómica de célula única: un perfil por célula}
  \begin{columns}[c]
    \begin{column}{0.5\linewidth}
      \centering
      \begin{tikzpicture}[font=\scriptsize, every node/.style={align=center}]
        % bulk
        \node[font=\scriptsize\bfseries] at (1.0,2.6) {bulk};
        \foreach \c in {0,...,3} { \fill[ucoblue!\the\numexpr 35+\c*12\relax,draw=white] (\c*0.5,2.1) rectangle ++(0.45,-0.35); }
        \node[font=\tiny,text=greyink] at (1.0,1.45) {1 muestra = 1 fila: promedio de todas las células};
        % single cell
        \node[font=\scriptsize\bfseries] at (1.0,0.85) {célula única};
        \foreach \r in {0,...,5} {
          \foreach \c in {0,...,3} {
            \pgfmathsetmacro{\v}{ifthenelse(rnd>0.72,rnd*65+35,0)}
            \fill[bluemid!\v!white,draw=greyink!25] (\c*0.5,0.4-\r*0.28) rectangle ++(0.45,-0.25);
          }
        }
        \node[font=\tiny,text=greyink] at (1.0,-1.55) {1 célula = 1 fila: miles de filas, casi todo ceros};
      \end{tikzpicture}
    \end{column}
    \begin{column}{0.48\linewidth}\small
      \begin{tabular}{@{}>{\raggedright\arraybackslash}p{0.2\linewidth}>{\raggedright\arraybackslash}p{0.25\linewidth}>{\raggedright\arraybackslash}p{0.45\linewidth}@{}}
        \toprule
        & \textbf{bulk} & \textbf{célula única}\\\midrule
        unidad & muestra & célula\\
        filas & decenas & $10^4$--$10^7$\\
        ceros & pocos & casi todas las entradas\\
        revela & el promedio & tipos celulares y células raras\\
        \bottomrule
      \end{tabular}
      \vspace{0.25cm}
      \begin{itemize}\footnotesize\setlength\itemsep{3pt}
        \item Cada célula recibe un \alert{código de barras}, de modo que las lecturas se asignan a su célula.
        \item Los tipos celulares se identifican después por clustering y \emph{genes marcadores} (p.\,ej.\ \gene{LYZ} en monocitos, \gene{CD3E} en linfocitos T).
      \end{itemize}
    \end{column}
  \end{columns}
\end{frame}

\begin{frame}{Célula única por indexado combinatorio (split-pool)}
  \centering
  \begin{tikzpicture}[scale=0.92,transform shape,font=\scriptsize, every node/.style={align=center}]
    % plates: 3 small grids
    \foreach \p/\x/\col in {1/0.3/ucoblue, 2/3.5/orangeslot, 3/6.7/greenslot} {
      \foreach \r in {0,...,2} \foreach \c in {0,...,3} {
        \pgfmathsetmacro{\sh}{20+mod((\r*4+\c)*7+\p*11,70)}
        \fill[\col!\sh!white] (\x+\c*0.3,-\r*0.3) rectangle ++(0.26,-0.26);
      }
      \node[font=\scriptsize\bfseries,text=\col] at (\x+0.6,0.3) {ronda \p};
    }
    \node[draw=greyink!50,rounded corners,minimum width=1.9cm,minimum height=0.8cm,font=\scriptsize] (cells) at (-1.5,-0.45) {células fijadas\\ \tiny (cada célula es su ``tubo'')};
    \draw[-{Stealth},thick] (cells) -- (0.1,-0.45);
    \draw[-{Stealth},thick] (1.65,-0.45) -- node[above,font=\tiny] {mezclar y repartir} (3.3,-0.45);
    \draw[-{Stealth},thick] (4.85,-0.45) -- node[above,font=\tiny] {mezclar y repartir} (6.5,-0.45);
    \node[draw=accentred,fill=accentred!8,rounded corners,minimum width=2.5cm,minimum height=0.95cm] (res) at (10.2,-0.45) {cada célula: un código\\ único \textbf{B1\,--\,B2\,--\,B3}};
    \draw[-{Stealth},thick] (8.05,-0.45) -- (res);
  \end{tikzpicture}
  \vspace{0.3cm}

  \begin{columns}[T]
    \begin{column}{0.54\linewidth}\small
      \begin{itemize}\setlength\itemsep{3pt}
        \item En cada ronda las células se \alert{reparten en pocillos al azar}; cada pocillo añade su propio código de barras al RNA de las células que contiene.
        \item Tras tres rondas, la probabilidad de que dos células compartan los mismos tres códigos es minúscula: con 96 códigos por ronda, $96^3\approx 880\,000$ combinaciones (ilustrativo).
        \item \textbf{Sin gotas ni aislamiento de células}: millones de células en un solo experimento.
      \end{itemize}
    \end{column}
    \begin{column}{0.42\linewidth}
      \begin{block}{En nuestros datos}
        \footnotesize
        El \alert{código de la primera ronda} es el pocillo de la placa de estimulación, es decir, la \alert{condición}: donante $\times$ citoquina. Es la columna \code{bc1\_well} del archivo crudo.\\[2pt]
        El control PBS se repartió en \alert{6 pocillos}.
      \end{block}
    \end{column}
  \end{columns}
\end{frame}

\begin{frame}{¿Por qué hay que hacer pseudobulk primero?}
  \centering
  \begin{tikzpicture}[font=\scriptsize]
    \foreach \r in {0,...,5} {
      \foreach \c in {0,...,7} {
        \pgfmathsetmacro{\v}{ifthenelse(rnd>0.5,rnd*65+35,0)}
        \fill[bluemid!\v!white,draw=greyink!25] (\c*0.28,-\r*0.28) rectangle ++(0.26,-0.26);
      }
    }
    \node[align=center] at (1.1,0.6) {células de un donante, una citoquina, un tipo celular\\ \textbf{células $\times$ genes} (conteos UMI)};
    \draw[decorate,decoration={brace,amplitude=4pt,mirror}] (-0.1,0) -- (-0.1,-1.65) node[midway,left=5pt] {células};
    \draw[-{Stealth},very thick,accentred] (2.5,-0.8) -- node[above] {$\sum$ sobre células} (4.5,-0.8);
    \foreach \c in {0,...,7} {
      \pgfmathsetmacro{\v}{rnd*45+50}
      \fill[ucoblue!\v] (4.8+\c*0.28,-0.55) rectangle ++(0.26,-0.26);
    }
    \node[align=center] at (5.9,0.2) {\textbf{una fila}\\ un perfil ``tipo bulk''};
    \draw[thick] (8.6,0.2) rectangle (11.6,-1.9);
    \foreach \r in {1,...,6} \draw[gray!50] (8.6,0.2-0.3*\r) -- (11.6,0.2-0.3*\r);
    \fill[ucoblue!60] (8.6,-0.7) rectangle (11.6,-1.0);
    \node[right] at (11.7,-0.85) {1 muestra};
    \node at (10.1,-2.3) {\textbf{X}: \numnsamples\ muestras $\times$ \numngenesraw\ genes};
  \end{tikzpicture}
  \vspace{0cm}

  \begin{itemize}\small\setlength\itemsep{2pt}
    \item \alert{La unidad de replicación es el donante, no la célula}: millones de células proceden de solo 12 personas, y las células de una persona no son independientes. Analizar las células como si fueran réplicas produce muchos falsos positivos (\emph{pseudorreplicación}).
    \item Los conteos de célula única son dispersos y ruidosos; sumados sobre miles de células se comportan como \alert{RNA-seq bulk}, así que se aplica una estadística estándar y transparente.
    \item Hace los datos pequeños: 9.7 millones de células $\to$ \numnsamples\ muestras.
  \end{itemize}
\end{frame}

% ────────────────────────────────────────────────────────────────────────────────────────
\begin{frame}{El conjunto de datos: el Diccionario Humano de Citoquinas}
  \begin{columns}[c]
    \begin{column}{0.52\linewidth}
      \centering
      \includegraphics[width=\linewidth,height=0.74\textheight,keepaspectratio]{paper_fig1ab.png}\\[-1pt]
      {\tiny Oesinghaus \emph{et al.}\ 2025, Fig.\ 1a--b (CC BY-NC-ND 4.0)}
    \end{column}
    \begin{column}{0.44\linewidth}\small
      \begin{itemize}\setlength\itemsep{4pt}
        \item \alert{12 donantes sanos} (6 mujeres, 6 hombres; 34--75 años)
        \item \alert{90 citoquinas} + un control PBS, in vitro durante \alert{24 h}
        \item $\sim$\alert{9.7 millones} de PBMC tras el control de calidad
        \item 16 tipos celulares, anotados con genes marcadores (panel b)
      \end{itemize}
      \vspace{0.2cm}
      \footnotesize
      \textbf{Pregunta del artículo:} ¿qué hace cada citoquina a cada tipo celular inmune y cuánto difieren las personas?
    \end{column}
  \end{columns}
\end{frame}

\begin{frame}{Cómo se generaron los datos}
  \centering
  \begin{tikzpicture}[font=\scriptsize, every node/.style={align=center},
      stp/.style={draw=ucoblue,fill=ucoblue!8,rounded corners,minimum height=1.45cm,minimum width=1.85cm,inner sep=2pt},
      numdot/.style={circle,fill=accentred,text=white,inner sep=1.5pt,font=\tiny\bfseries}]
    \node[stp] (a) {PBMC de\\ 12 donantes};
    \node[stp,right=0.3cm of a] (b) {placa de 96\\ 10$^6$ células/pocillo\\ citoquina o PBS\\ \textbf{24 h}};
    \node[stp,right=0.3cm of b] (c) {fijar las\\ células};
    \node[stp,right=0.3cm of c] (d) {\textbf{3 rondas}\\ de barcoding\\ split-pool};
    \node[stp,right=0.3cm of d] (e) {secuenciación\\ $\sim$31\,000\\ lecturas/célula};
    \node[stp,right=0.3cm of e,fill=accentred!10,draw=accentred] (f) {conteos\\ célula $\times$ gen};
    \foreach \x/\y in {a/b,b/c,c/d,d/e,e/f} \draw[-{Stealth},thick] (\x) -- (\y);
    \node[stp,below=0.75cm of d,fill=white] (g) {control de calidad\\ eliminar dobletes};
    \node[stp,left=0.3cm of g,fill=white] (h) {clustering y\\ anotación:\\ 16 tipos celulares};
    \node[stp,right=0.3cm of g,fill=white,draw=accentred] (i) {\textbf{pseudobulk}\\ donante $\times$ citoquina\\ $\times$ tipo celular};
    \draw[-{Stealth},thick] (f.south) |- ($(g.north)+(0.5,0.35)$) -- ($(g.north)+(0.5,0)$);
    \draw[-{Stealth},thick] (g) -- (i); \draw[-{Stealth},thick] (h) -- (g);
  \end{tikzpicture}
  \vspace{0.25cm}

  \footnotesize
  \begin{itemize}\setlength\itemsep{2pt}
    \item Barcoding con el protocolo GigaLab de Parse Biosciences; sublibrerías de $\sim$100\,000 células; secuenciación en un Ultima UG-100.
    \item Se conservaron células con 400--7\,000 genes detectados y $<15\,\%$ de lecturas mitocondriales.
  \end{itemize}
\end{frame}

\begin{frame}{¿Qué contiene el subconjunto con el que trabajamos?}
  \begin{columns}[T]
    \begin{column}{0.56\linewidth}
      \small
      \begin{tabular}{@{}>{\raggedright\arraybackslash}p{0.2\linewidth}>{\raggedright\arraybackslash}p{0.22\linewidth}>{\raggedright\arraybackslash}p{0.52\linewidth}@{}}
        \toprule
        & \textbf{artículo} & \textbf{nuestra práctica}\\\midrule
        donantes & 12 & 12\\[1pt]
        citoquinas & 90 + PBS & 87 + PBS\\[1pt]
        tipos celulares & 16 & 4: T CD4, T CD8, NK CD56\textsuperscript{low}, monocitos CD14\\[1pt]
        unidad & células únicas & muestras \alert{pseudobulk}\\[1pt]
        genes & todos & \numngenesraw\ codificantes de proteínas ($\to$ \numngenes\ tras filtrar)\\
        \bottomrule
      \end{tabular}
    \end{column}
    \begin{column}{0.4\linewidth}
      \begin{block}{El diseño}
        \centering
        $4$ tipos celulares\\ $\times$ $88$ condiciones\\ \footnotesize(87 citoquinas + PBS)\\ \normalsize $\times$ $12$ donantes\\[3pt]
        $=$ \alert{\numnsamples\ muestras}
      \end{block}
      \footnotesize
      El control PBS se midió en 6 pocillos; el archivo guarda cada pocillo \emph{y} su total. Nos quedamos con el total (extra B1 del notebook).
    \end{column}
  \end{columns}
\end{frame}

\begin{frame}{AnnData: una estructura de datos de Python para datos ómicos}
  \begin{columns}[c]
    \begin{column}{0.38\linewidth}
      \centering
      \includegraphics[width=\linewidth,height=0.76\textheight,keepaspectratio]{anndata_schema.png}\\[-1pt]
      {\tiny Esquema de AnnData (documentación de anndata)}
    \end{column}
    \begin{column}{0.58\linewidth}\small
      \begin{itemize}\setlength\itemsep{3pt}
        \item \textbf{Un objeto con varias tablas que siempre permanecen alineadas.}
        \item \textcolor{ucoblue}{\textbf{\code{X}}}: la matriz principal, $n_{\text{obs}}\times n_{\text{vars}}$
        \item \textcolor{yellowslot}{\textbf{\code{obs}}}: una fila por \alert{observación} (muestra o célula): metadatos
        \item \textcolor{bluemid}{\textbf{\code{var}}}: una fila por \alert{variable} (gen): metadatos
        \item \textcolor{greenslot}{\textbf{\code{layers}}}: más matrices, con la misma forma que \code{X}
        \item \code{obsm}, \code{varm}, \code{obsp}, \code{varp}, \code{uns}: embeddings (PCA), grafos, resultados de formato libre
      \end{itemize}
      \vspace{0.15cm}
      \begin{block}{Por qué importa}
        \footnotesize Seleccionar un subconjunto mueve \emph{todas} las tablas a la vez: ningún gen ni muestra puede desalinearse de sus metadatos. Se guarda como archivo \code{.h5ad}.
      \end{block}
    \end{column}
  \end{columns}
\end{frame}

\begin{frame}[fragile]{AnnData en nuestra práctica}
  \begin{columns}[T]
    \begin{column}{0.52\linewidth}
      \small
      \begin{tabular}{@{}>{\ttfamily}p{0.16\linewidth}>{\raggedright\arraybackslash}p{0.47\linewidth}>{\raggedright\arraybackslash}p{0.3\linewidth}@{}}
        \toprule
        \textrm{\textbf{slot}} & \textbf{nuestros datos} & \textbf{forma}\\\midrule
        X & conteos crudos, luego $\log_2$(CPM+1) & \numnsamples\,$\times$\,\numngenesraw\\[2pt]
        obs & donor, cytokine, treatment, cell\_type, número de células & \numnsamples\ filas\\[2pt]
        var & nombres de genes, conteos totales & \numngenesraw\ filas\\[2pt]
        layers & \code{counts}, \code{cpm} & como \code{X}\\[2pt]
        obsm & puntuaciones del PCA (Parte 3) & \numnsamples\,$\times$\,10\\
        \bottomrule
      \end{tabular}
    \end{column}
    \begin{column}{0.46\linewidth}
\begin{lstlisting}
# primero filas (muestras), luego genes
adata[["Donor1|CD14_Mono|IFN-gamma"],
      ["STAT1"]]

mask = adata.obs["cell_type"] == "CD14_Mono"
adata[mask, ["IDO1", "CXCL10"]]
\end{lstlisting}
      \footnotesize La primera selección es de $1\times1$; la segunda, de $1\,056\times2$ (todas las muestras de monocitos).
    \end{column}
  \end{columns}
  \vspace{0.25cm}
  \begin{block}{Manos a la obra}
    \small Volvemos al notebook: \textbf{preparación de la Parte 1} (instalar, cargar los datos) y después la \textbf{Parte 2}.
    \timebox{unos 25 min para la Parte 1}
  \end{block}
\end{frame}

% ════════════════════════════════════════════════════════════════════════════════════════
\section{Parte 2: Manipular y normalizar los datos}

\begin{frame}{Hoja de ruta de la Parte 2: 40 min, en el notebook}
  \small
  \begin{tabular}{@{}p{0.07\linewidth}p{0.6\linewidth}p{0.2\linewidth}@{}}
    \toprule
    \textbf{Paso} & \textbf{Pregunta} & \textbf{Tiempo}\\\midrule
    \textbf{2.1} & Preparar el objeto & 3 min\\[3pt]
    \textbf{2.2} & ¿Qué pinta tienen los datos? De \textbf{1 número} a la \textbf{matriz completa} & 10 min\\[3pt]
    \textbf{2.3} & ¿Qué es un \textbf{conteo}? ¿Por qué las librerías tienen distinto tamaño? & 7 min\\[3pt]
    \textbf{2.4} & \textbf{Normalización}: qué, por qué, cómo & 10 min\\[3pt]
    \textbf{2.5} & ¿Por qué \textbf{filtramos} los genes poco expresados? & 5 min\\[3pt]
    \textbf{2.6} & \textbf{Mapas de calor} antes y después de normalizar & 5 min\\
    \bottomrule
  \end{tabular}
  \vspace{0.3cm}

  \begin{block}{Objetivo}
    Entender qué hay dentro de los datos y convertir los \alert{conteos crudos} en números que podamos \alert{comparar} entre muestras y entre genes.
  \end{block}
\end{frame}

\begin{frame}{2.2 De un número a la matriz completa}
  \centering
  \begin{tikzpicture}[font=\scriptsize]
    % the big matrix
    \draw[thick,fill=softgrey] (0,0) rectangle (6.2,-3.2);
    \node[align=center] at (3.4,-1.9) {\textbf{la matriz completa}\\ \numnsamples\ muestras $\times$ \numngenesraw\ genes\\ \footnotesize 76 millones de números};
    \node[above,font=\tiny,text=greyink] at (3.1,0) {genes $\rightarrow$};
    \node[rotate=90,above,font=\tiny,text=greyink] at (0,-1.6) {muestras $\rightarrow$};
    % the three zoom levels, at the corner
    \fill[accentred] (0,0) rectangle ++(0.2,-0.2);
    \fill[bluemid]   (0,-0.3) rectangle ++(1.0,-0.2);
    \fill[orangeslot] (0,-0.65) rectangle ++(1.0,-1.0);
    \foreach \i in {1,...,4} { \draw[white] (\i*0.2,-0.65) -- (\i*0.2,-1.65); \draw[white] (0,-0.65-\i*0.2) -- (1.0,-0.65-\i*0.2); }
    % explanations
    \node[anchor=west,draw=accentred,very thick,rounded corners,fill=accentred!6,minimum width=6.0cm,align=left,inner sep=3pt] (a) at (7.3,0.1)
      {\textbf{$1\times1$} un gen en una muestra\\ \ttfamily\scriptsize adata[[sample], ["STAT1"]]};
    \node[anchor=west,draw=bluemid,very thick,rounded corners,fill=bluemid!6,minimum width=6.0cm,align=left,inner sep=3pt] (b) at (7.3,-1.15)
      {\textbf{$1\times5$} una muestra, cinco genes\\ \ttfamily\scriptsize adata[[sample], five\_genes]};
    \node[anchor=west,draw=orangeslot,very thick,rounded corners,fill=orangeslot!6,minimum width=6.0cm,align=left,inner sep=3pt] (c) at (7.3,-2.4)
      {\textbf{$5\times5$} cinco muestras, cinco genes\\ \ttfamily\scriptsize adata[five\_samples, five\_genes]};
  \end{tikzpicture}

  \vspace{0.2cm}
  \begin{itemize}\small\setlength\itemsep{3pt}
    \item Cada trozo viene con sus \alert{metadatos}: \code{.obs} describe las muestras (filas) y \code{.var} los genes (columnas).
    \item La tabla de 5\,$\times$\,5 es solo la esquina superior izquierda de la matriz completa: lo que hagamos a continuación se aplica a \emph{toda} ella.
  \end{itemize}
\end{frame}

\begin{frame}{2.3 ¿Qué es un conteo?}
  \centering
  \begin{tikzpicture}[font=\scriptsize, every node/.style={align=center}]
    % molecules
    \node[font=\scriptsize\bfseries] at (1.2,1.15) {3 moléculas de mRNA de\\ un gen (todas las células sumadas)};
    \foreach \i/\col/\u in {0/accentred/AAGT, 1/bluemid/CCTA, 2/greenslot/GGAC} {
      \draw[thick,\col] (\i*0.95,0.1) .. controls (\i*0.95+0.2,0.4) and (\i*0.95+0.4,-0.2) .. (\i*0.95+0.7,0.1);
      \node[fill=\col,text=white,inner sep=1.5pt,font=\tiny\ttfamily] at (\i*0.95+0.35,-0.28) {\u};
    }
    \draw[-{Stealth},very thick] (3.2,0) -- node[above,font=\tiny] {amplificación\\ + secuenciación} (5.2,0);
    % reads
    \node[font=\scriptsize\bfseries] at (7.0,1.15) {6 lecturas};
    \foreach \i/\col/\u in {0/accentred/AAGT, 1/accentred/AAGT, 2/bluemid/CCTA, 3/greenslot/GGAC, 4/greenslot/GGAC, 5/greenslot/GGAC} {
      \node[fill=\col,text=white,inner sep=2pt,font=\tiny\ttfamily] at (5.6+\i*0.62,0) {\u};
    }
    \draw[-{Stealth},very thick] (9.6,0) -- node[above,font=\tiny] {colapsar lecturas\\ con el mismo UMI} (11.6,0);
    \node[draw=accentred,rounded corners,fill=accentred!8,inner sep=4pt,font=\small\bfseries] at (12.9,0) {conteo = 3};
  \end{tikzpicture}
  \vspace{0.3cm}

  \begin{columns}[T]
    \begin{column}{0.52\linewidth}\small
      \begin{itemize}\setlength\itemsep{3pt}
        \item Un \textbf{UMI} (\emph{identificador molecular único}) es una etiqueta aleatoria añadida a cada molécula \emph{original}.
        \item Contar \alert{UMI distintos} en lugar de lecturas elimina los duplicados de PCR: un UMI $=$ una molécula.
      \end{itemize}
    \end{column}
    \begin{column}{0.44\linewidth}\small
      \begin{itemize}\setlength\itemsep{3pt}
        \item En un pseudobulk el conteo se \alert{suma sobre todas las células} de la muestra.
        \item Así, un conteo es un entero no negativo que depende de la \textbf{biología} y de \textbf{cuánto material} tenemos.
      \end{itemize}
    \end{column}
  \end{columns}
\end{frame}

\begin{frame}{2.3 Tamaño de librería: células $\times$ UMI por célula}
  \centering
  \figfit[0.42]{p2_libsize}\\[-1pt]
  \begin{columns}[c]
    \begin{column}{0.4\linewidth}
      \centering\scriptsize
      \begin{tabular}{lrrr}\toprule
 & células & UMI / célula & librería\\\midrule
CD4 T & 3\,041 & 2\,989 & 8.5\,M\\
CD8 T & 1\,302 & 3\,104 & 3.8\,M\\
NK & 407 & 3\,147 & 1.2\,M\\
Monocitos & 1\,238 & 6\,747 & 7.6\,M\\
\bottomrule\end{tabular}

    \end{column}
    \begin{column}{0.58\linewidth}\footnotesize
      \begin{itemize}\setlength\itemsep{2pt}
        \item \textbf{Tamaño de librería} $L_i=\sum_j\text{conteo}_{ij}$: conteos totales de la muestra.
        \item T CD4: los más abundantes; NK: los más escasos ($\sim$400 células).
        \item Monocitos: \alert{\numumiMono\ UMI por célula} frente a $\sim$\numumiNK\ en linfocitos (más mRNA).
        \item Controles PBS (suma de 6 pocillos): \alert{\numpbsratio$\times$ más profundos}.
      \end{itemize}
    \end{column}
  \end{columns}
\end{frame}

\begin{frame}{2.4 ¿Por qué normalizar? Un ejemplo}
  \begin{columns}[c]
    \begin{column}{0.56\linewidth}
      \figfit[0.62]{p2_toy}
    \end{column}
    \begin{column}{0.42\linewidth}\small
      \textit{¿Se usa más la palabra ``citoquina'' en un artículo de 10 páginas o en una revisión de 100?} Comparamos \alert{frecuencias}, no conteos.
      \vspace{0.2cm}
      \begin{itemize}\setlength\itemsep{3pt}
        \item Mismo donante: un control PBS (profundo) y una muestra con IFN-$\gamma$ (5$\times$ más superficial). El gen B es un gen de respuesta a interferón.
        \item \alert{Conteos crudos}: todos los genes parecen \emph{menores} con IFN-$\gamma$, incluso el gen B, inducido por interferón.
        \item \alert{CPM}: A, C y D no cambian y B sube $4\times$. La respuesta correcta.
      \end{itemize}
    \end{column}
  \end{columns}
\end{frame}

\begin{frame}{CPM: conteos por millón}
  \[
    \text{CPM}_{ij}\;=\;\frac{\text{conteo}_{ij}}{L_i}\times 10^{6},\qquad L_i=\sum_{j}\text{conteo}_{ij}
  \]
  \centering
  \figfit[0.46]{p2_actb}\\[-2pt]
  \small \gene{ACTB}, un gen constitutivo, en monocitos (incluidas las muestras PBS). Tras el CPM ya no sigue al tamaño de librería.
\end{frame}

\begin{frame}{El logaritmo: entre genes}
  \begin{columns}[c]
    \begin{column}{0.6\linewidth}
      \figfit[0.6]{p2_logvar}
    \end{column}
    \begin{column}{0.38\linewidth}\small
      \begin{itemize}\setlength\itemsep{3pt}
        \item En CPM, la \alert{varianza crece mucho con la media}.
        \item Los 20 genes más variables acumulan el \numtoptwentycpm\ de la varianza total en CPM\ldots{}
        \item \ldots{}y solo el \numtoptwentylog\ tras $\log_2(\text{CPM}+1)$.
        \item El PCA, el clustering y las pruebas $t$ dependen de la varianza: sin el logaritmo, unos pocos genes abundantes lo deciden todo.
      \end{itemize}
    \end{column}
  \end{columns}
\end{frame}

\begin{frame}{El logaritmo y el pseudoconteo}
  \[ \text{log-CPM}_{ij}=\log_2\bigl(\text{CPM}_{ij}+1\bigr) \]
  \begin{columns}[T]
    \begin{column}{0.45\linewidth}
      \footnotesize\textbf{Los fold changes se vuelven simétricos}\\[3pt]
      \scriptsize\setlength\tabcolsep{2.5pt}
      \begin{tabular}{@{}lccccccc@{}}
        \toprule
        fold change & $\times4$ & $\times2$ & $\times1.4$ & $\times1$ & $\times0.7$ & $\times0.5$ & $\times0.25$\\
        $\lfc$ & $+2$ & $+1$ & $+0.5$ & $0$ & $-0.5$ & $-1$ & $-2$\\
        \bottomrule
      \end{tabular}\\[4pt]
      La diferencia de dos valores en log \emph{es} el log fold change.
    \end{column}
    \begin{column}{0.53\linewidth}
      \footnotesize\textbf{¿Por qué ``+1''?}\\[3pt]
      \setlength\tabcolsep{2.5pt}
      \begin{tabular}{@{}lrrrr@{}}
        \toprule
        gen & ctrl & trt & $\lfc$ & $\lfc^{+1}$\\\midrule
        casi silencioso & 0.1 & 0.4 & $+2.00$ & $+0.35$\\
        silencioso $\to$ detectado & 0 & 0.5 & $+\infty$ & $+0.58$\\
        moderado & 20 & 40 & $+1.00$ & $+0.97$\\
        alto & 500 & 1000 & $+1.00$ & $+1.00$\\
        \bottomrule
      \end{tabular}\\[2pt]
      \scriptsize ctrl, trt = CPM; $\lfc^{+1}$ usa CPM\,$+1$.\\[3pt]
      El pseudoconteo contrae los fold changes de números minúsculos: ``no me fío de cocientes de una o dos moléculas''.
    \end{column}
  \end{columns}
\end{frame}

\begin{frame}{2.5 ¿Por qué filtrar los genes poco expresados?}
  \centering
  \figfit[0.38]{p2_poisson}\\[-2pt]
  \small El \alert{mismo gen medido dos veces} (sin cambio real), con cuatro niveles de expresión (solo ruido de Poisson)
  \begin{itemize}\small\setlength\itemsep{3pt}
    \item Los conteos bajos son \alert{puro ruido de muestreo}: con 2 frente a 4 moléculas, un ``doblaje'' es mala suerte.
    \item Dan \alert{fold changes falsos} y \alert{más pruebas} que corregir (Parte 4).
  \end{itemize}
\end{frame}

\begin{frame}{$\geq$ 10 conteos en $\geq$ 12 muestras}
  \begin{columns}[c]
    \begin{column}{0.56\linewidth}
      \figfit[0.6]{p2_filter}
    \end{column}
    \begin{column}{0.42\linewidth}\small
      \begin{itemize}\setlength\itemsep{4pt}
        \item 12 = el número de donantes: el tamaño de un grupo experimental (la idea de \texttt{filterByExpr} en edgeR).
        \item Un gen que nunca llega a 10 conteos en todo un grupo no se puede analizar en ningún sitio.
        \item Superan el filtro \alert{\numngenes\ de \numngenesraw\ genes}. Los \numngenesremoved\ eliminados acumulan el 0.011\,\% de los conteos (p.\,ej.\ receptores olfativos).
        \item Estos números son un consejo, no una ley: puedes cambiar el umbral en el notebook y ver qué pasa.
      \end{itemize}
    \end{column}
  \end{columns}
\end{frame}

\begin{frame}{2.6 Mapas de calor: marcadores de tipo celular, antes y después}
  \centering
  \figfit[0.72]{p2_heat_markers}
  \vspace{-0.15cm}\\
  \scriptsize
  \textbf{z-score por gen entre muestras:} $z_{gi}=(x_{gi}-\bar x_g)/s_g$ (rojo: por encima de la media del gen). Columnas: media de los 12 donantes. \alert{Antes}: las columnas PBS (las librerías más profundas) destacan. \alert{Después}: bloques limpios por tipo celular.
\end{frame}

\begin{frame}{2.6 Mapas de calor: genes de la respuesta inmune}
  \begin{columns}[c]
    \begin{column}{0.71\linewidth}
      \figfit[0.8]{p2_heat_immune}
    \end{column}
    \begin{column}{0.27\linewidth}\scriptsize
      Cuatro citoquinas, cuatro rutas:
      \begin{itemize}\setlength\itemsep{3pt}
        \item \textbf{IFN-$\gamma$ / STAT1}: \gene{GBP1}, \gene{IDO1}, \gene{CXCL9}\ldots, más fuerte en monocitos
        \item \textbf{IL-4 / STAT6}: \gene{ALOX15}, \gene{CCL17}, \gene{FCER2}, monocitos
        \item \textbf{IL-1$\beta$ / NF-$\kappa$B}: \gene{IL6}, \gene{CXCL8}, \gene{CCL2}, monocitos
        \item \textbf{IL-2 / STAT5}: \gene{IL2RA}, \gene{GZMB}, \gene{CISH}, NK y linfocitos T
      \end{itemize}
      \alert{Normalizado: biología. Crudo: profundidad de secuenciación.}\\[2pt]
      Los z-scores son solo para visualizar.
    \end{column}
  \end{columns}
\end{frame}

\begin{frame}{Tu turno en el notebook (Parte 2)}
  \small
  Las celdas que empiezan por \code{\# CAMBIA ESTO} contienen parámetros que puedes modificar. Cámbialos, vuelve a ejecutar y mira qué pasa:
  \vspace{0.2cm}

  \begin{tabular}{@{}>{\ttfamily}p{0.27\linewidth}>{\raggedright\arraybackslash}p{0.68\linewidth}@{}}
    \toprule
    \textrm{\textbf{Parámetro}} & \textbf{Prueba\ldots}\\\midrule
    sample, genes5, samples5 & otras muestras y genes en las tablas de $1\times1$ y $5\times5$: ¿qué marcador es alto en qué tipo celular?\\[2pt]
    mask & otro tipo celular y otras citoquinas: ¿es mayor el conteo crudo con la citoquina?\\[2pt]
    MIN\_COUNT, MIN\_SAMPLES & un filtro de genes más estricto o más laxo: ¿cuántos genes sobreviven?\\[2pt]
    CELL\_TYPE\_MARKERS, IMMUNE\_GENES & tus propios genes en los mapas de calor antes/después\\
    \bottomrule
  \end{tabular}
\end{frame}

\begin{frame}{La Parte 2 en una diapositiva}
  \footnotesize
  \begin{tabular}{@{}c>{\raggedright\arraybackslash}p{0.14\linewidth}>{\raggedright\arraybackslash}p{0.33\linewidth}>{\raggedright\arraybackslash}p{0.22\linewidth}>{\raggedright\arraybackslash}p{0.17\linewidth}@{}}
    \toprule
    & \textbf{Paso} & \textbf{Corrige} & \textbf{Fórmula} & \textbf{Se guarda en}\\\midrule
    1 & CPM & distintos tamaños de librería (entre muestras) & $\text{conteo}_{ij}/L_i\times10^6$ & \code{layers["cpm"]}\\[2pt]
    2 & $\log_2(\cdot+1)$ & rango dinámico, relación media--varianza (entre genes) & $\log_2(\text{CPM}+1)$ & \code{adata.X}\\[2pt]
    3 & filtro de genes & ruido y exceso de pruebas & $\geq10$ conteos en $\geq12$ muestras & filas de \code{adata}\\[2pt]
    4 & z-score por gen & la línea base del gen oculta el patrón (mapas de calor) & $(x_{gi}-\bar x_g)/s_g$ & no se guarda\\
    \bottomrule
  \end{tabular}
  \vspace{0.4cm}
  \begin{block}{Recuerda}
    \code{layers["counts"]} = conteos crudos \quad $\to$ \quad \code{layers["cpm"]} = CPM \quad $\to$ \quad \code{adata.X} = log$_2$(CPM+1)
  \end{block}
\end{frame}

\begin{frame}[standout]
  Descanso\\[6pt]
  {\Large 10 minutos}
\end{frame}

% ════════════════════════════════════════════════════════════════════════════════════════
\section{Parte 3: Análisis de Componentes Principales}

\begin{frame}{Hoja de ruta de la Parte 3: un análisis para gobernarlos a todos (35 min)}
  \footnotesize
  \begin{columns}[T]
    \begin{column}{0.54\linewidth}
      \textbf{Diapositivas} (10 min)
      \begin{itemize}\setlength\itemsep{1pt}
        \item ¿Por qué reducir la dimensionalidad?
        \item El PCA como \alert{rotación}
        \item \alert{SVD}: cómo se calcula
        \item \alert{Proyección}: puntuaciones y cargas
        \item El PCA en distintos campos
      \end{itemize}
      \vspace{0.1cm}
      \textbf{Notebook} (25 min)
      \begin{itemize}\setlength\itemsep{1pt}
        \item 3.1 PCA a mano con 2 genes
        \item 3.2 PCA de todas las muestras
        \item 3.3 cuatro coloreados: tipo celular, donante, citoquina, control vs tratado
        \item 3.4 ¿qué significa cada PC? R$^2$ y cargas
        \item 3.5 acercamiento a un tipo celular
      \end{itemize}
    \end{column}
    \begin{column}{0.42\linewidth}
      \begin{block}{La pregunta}
        Cada muestra se describe con $\sim$15\,500 números (uno por gen). \textbf{¿Qué impulsa la variación entre muestras?}
      \end{block}
      \begin{block}{La idea}
        Encontrar unos pocos ejes nuevos que conserven la mayor variación posible y después \alert{colorear las muestras} para ver qué significan.
      \end{block}
    \end{column}
  \end{columns}
\end{frame}

\begin{frame}{¿Por qué reducir la dimensionalidad?}
  \begin{columns}[c]
    \begin{column}{0.5\linewidth}\small
      \begin{itemize}\setlength\itemsep{4pt}
        \item \numnsamples\ muestras $\times$ \numngenes\ genes: imposible de mirar.
        \item Muchos genes se mueven \alert{juntos} (la misma citoquina enciende todo un programa): los datos son más simples de lo que sugieren sus 15\,000 dimensiones.
        \item El \textbf{PCA} busca unas pocas \alert{combinaciones de genes} que resuman esa variación compartida:
        \begin{itemize}\footnotesize
          \item \alert{lineal}: sumas ponderadas de genes
          \item \alert{ordenado}: la PC1 tiene la mayor varianza, la PC2 la mayor de lo que queda, \ldots
          \item \alert{no supervisado}: no sabe nada de citoquinas ni de donantes
        \end{itemize}
      \end{itemize}
    \end{column}
    \begin{column}{0.46\linewidth}
      \centering
      \begin{tikzpicture}[font=\scriptsize]
        % a 3D-ish object and its shadow
        \draw[thick,fill=bluemid!15] (0,0) -- (1.6,0.3) -- (2.3,1.4) -- (0.7,1.1) -- cycle;
        \draw[thick,fill=bluemid!30] (0.7,1.1) -- (2.3,1.4) -- (2.0,2.1) -- (0.5,1.8) -- cycle;
        \node at (1.2,-0.35) {muchas dimensiones};
        \draw[-{Stealth},very thick,accentred] (2.9,0.9) -- node[above] {proyectar} (4.2,0.9);
        \draw[thick,fill=softgrey] (4.6,0.2) rectangle (6.6,1.7);
        \foreach \x/\y in {4.9/0.5,5.2/1.2,5.5/0.7,5.8/1.4,6.1/0.9,6.3/0.4,5.0/0.95} \fill[bluemid] (\x,\y) circle (1.5pt);
        \node at (5.6,-0.35) {2 dimensiones};
      \end{tikzpicture}\\[4pt]
      \footnotesize Una buena proyección conserva la \alert{estructura}: las muestras parecidas siguen cerca.
    \end{column}
  \end{columns}
\end{frame}

\begin{frame}{El PCA es una rotación}
  \begin{columns}[c]
    \begin{column}{0.52\linewidth}
      \centering
      \begin{tikzpicture}[font=\scriptsize]
        \pgfmathsetseed{11}
        % original axes
        \draw[-{Stealth},thick,greyink!70] (-3.1,0) -- (3.3,0) node[right] {gen A};
        \draw[-{Stealth},thick,greyink!70] (0,-2.1) -- (0,2.3) node[above] {gen B};
        \begin{scope}[rotate=35]
          \foreach \i in {1,...,170} {
            \pgfmathsetmacro{\px}{(rnd+rnd+rnd-1.5)*2.0}
            \pgfmathsetmacro{\py}{(rnd+rnd+rnd-1.5)*0.7}
            \fill[bluemid!75] (\px,\py) circle (1.4pt);
          }
          \draw[-{Stealth},very thick,accentred] (-2.6,0) -- (2.8,0) node[right,font=\footnotesize\bfseries] {PC1};
          \draw[-{Stealth},very thick,accentred] (0,-1.3) -- (0,1.4) node[above,font=\footnotesize\bfseries] {PC2};
        \end{scope}
      \end{tikzpicture}
    \end{column}
    \begin{column}{0.46\linewidth}\small
      \begin{enumerate}\setlength\itemsep{4pt}
        \item \textbf{Centrar} cada gen (restarle su media): el origen se traslada al centro de la nube.
        \item \textbf{Rotar} los ejes hasta que el primero apunte hacia donde la nube es más larga: esa es la \alert{PC1}. La PC2 es perpendicular, y así sucesivamente.
        \item No se pierde nada: una rotación conserva todas las distancias entre muestras.
      \end{enumerate}
      \vspace{0.2cm}
      \begin{block}{Reducción de dimensionalidad}
        \footnotesize = quedarse solo con las primeras PCs y olvidar el resto.
      \end{block}
    \end{column}
  \end{columns}
\end{frame}

\begin{frame}{¿Cómo se calcula? La SVD}
  \centering
  \begin{tikzpicture}[font=\scriptsize, every node/.style={align=center}]
    \node[draw=greyink,fill=softgrey,minimum width=1.7cm,minimum height=2.3cm] (X) at (0,0) {};
    \node at (0,0) {\Large $X_c$};
    \node at (1.15,0) {\Large $=$};
    \node[draw=orangeslot,fill=orangeslot!15,minimum width=1.0cm,minimum height=2.3cm] (U) at (2.2,0) {};
    \node at (2.2,0) {\Large $U$};
    \node[draw=greyink,fill=greyink!20,minimum width=1.0cm,minimum height=1.0cm] (S) at (3.55,0.65) {};
    \node at (3.55,0.65) {\Large $S$};
    \draw[greyink] (3.1,1.1) -- (4.0,0.2);
    \node[draw=bluemid,fill=bluemid!15,minimum width=1.7cm,minimum height=1.0cm] (V) at (5.3,0.65) {};
    \node at (5.3,0.65) {\Large $V^{\!\top}$};
    \node[text=greyink,below=0.1cm of X] {datos centrados\\ muestras $\times$ genes};
    \node[text=orangeslot,below=0.1cm of U] {muestras\\ $\times$ PCs};
    \node[text=greyink,below=0.7cm of S] {fuerza de\\ cada PC};
    \node[text=bluemid,below=0.7cm of V] {PCs $\times$ genes\\ \alert{(cargas)}};
  \end{tikzpicture}
  \vspace{0.3cm}

  \begin{columns}[T]
    \begin{column}{0.5\linewidth}\small
      \begin{itemize}\setlength\itemsep{3pt}
        \item $V$ es una \alert{rotación}: sus columnas son los nuevos ejes (longitud unidad, perpendiculares): los pesos de cada gen en cada PC.
        \item $S$: cuánto se estiran los datos a lo largo de cada eje.
      \end{itemize}
    \end{column}
    \begin{column}{0.46\linewidth}\small
      \begin{itemize}\setlength\itemsep{3pt}
        \item \textbf{Varianza explicada} por la PC $k$: $\;s_k^2\big/\sum_j s_j^2$
        \item En Python: \code{np.linalg.svd(Xc)}; el PCA de \code{scikit-learn} hace exactamente esto.
      \end{itemize}
    \end{column}
  \end{columns}
\end{frame}

\begin{frame}{Proyección: puntuaciones y cargas}
  \centering
  \figfit[0.5]{p3_twogenes}
  \vspace{0.05cm}
  \begin{columns}[T]
    \begin{column}{0.5\linewidth}\small
      \begin{itemize}\setlength\itemsep{2pt}
        \item Monocitos, dos genes (\gene{ISG15}, \gene{MX1}).
        \item Las \alert{puntuaciones} (scores), es decir, las coordenadas de las muestras en los nuevos ejes, son su \textbf{proyección}: $T=X_c\,V=U\,S$
      \end{itemize}
    \end{column}
    \begin{column}{0.46\linewidth}\small
      \begin{itemize}\setlength\itemsep{2pt}
        \item Las dos cargas de la PC1 son positivas: PC1 $=$ ``nivel global de genes de interferón''.
        \item Solo la PC1 acumula el \alert{98\,\%} de la varianza.
      \end{itemize}
    \end{column}
  \end{columns}
\end{frame}

\begin{frame}[fragile]{El PCA en la práctica}
  \centering
  \begin{tikzpicture}[font=\scriptsize, every node/.style={align=center},
      stp/.style={draw=ucoblue,fill=ucoblue!8,rounded corners,minimum height=1.15cm,minimum width=2.15cm,inner sep=2pt}]
    \node[stp] (a) {datos normalizados\\ $\log_2$(CPM+1)};
    \node[stp,right=0.3cm of a] (b) {quedarse con los\\ \textbf{2\,000 genes más}\\ \textbf{variables}};
    \node[stp,right=0.3cm of b] (c) {centrar\\ los genes};
    \node[stp,right=0.3cm of c] (d) {\textbf{SVD}};
    \node[stp,right=0.3cm of d,fill=accentred!10,draw=accentred] (e) {puntuaciones, cargas,\\ varianza explicada};
    \foreach \x/\y in {a/b,b/c,c/d,d/e} \draw[-{Stealth},thick] (\x) -- (\y);
  \end{tikzpicture}
  \vspace{0.25cm}

\begin{lstlisting}
Mc = M - M.mean(axis=0)                          # 1) centrar los genes
U, S, Vt = np.linalg.svd(Mc, full_matrices=False)   # 2) SVD
scores   = U * S                                 # 3) proyeccion: muestras en las PCs
loadings = Vt.T                                  #    pesos de los genes
var_pct  = S**2 / (S**2).sum() * 100             #    % de varianza de cada PC
\end{lstlisting}
  \vspace{0.1cm}
  \small
  Después \alert{coloreamos las muestras} según los metadatos para averiguar qué significa cada PC.
\end{frame}

\begin{frame}{El PCA en distintos campos}
  \centering
  \begin{tikzpicture}[font=\scriptsize,
      card/.style={draw=ucoblue!50,rounded corners,fill=ucoblue!5,text width=3.95cm,inner sep=4pt,align=left,minimum height=2.15cm,anchor=north west}]
    \node[card] at (0,0)      {\textbf{\normalsize Genética de poblaciones}\\[2pt] Genotipos de miles de personas: PC1 y PC2 redibujan el mapa de la ascendencia y corrigen la estructura poblacional en estudios de asociación.};
    \node[card] at (4.2,0)    {\textbf{\normalsize Célula única}\\[2pt] Normalizar, genes variables, PCA: las primeras 30--50 PCs alimentan el grafo de vecinos, el clustering y el UMAP.};
    \node[card] at (8.4,0)    {\textbf{\normalsize QC de RNA-seq bulk}\\[2pt] Colorear por lote, condición o profundidad: aparecen atípicos, cambios de muestra y efectos de lote antes de cualquier prueba.};
    \node[card] at (0,-2.45)  {\textbf{\normalsize Otras ómicas}\\[2pt] Proteómica, metabolómica, microbioma: cientos de variables correlacionadas en pocos ejes; las cargas dicen cuáles los impulsan.};
    \node[card] at (4.2,-2.45) {\textbf{\normalsize Imágenes y señales}\\[2pt] ``Eigenfaces'', compresión de imágenes o espectros: se conservan unas pocas PCs.};
    \node[card] at (8.4,-2.45) {\textbf{\normalsize Cualquier tabla ancha}\\[2pt] Sensores, química, finanzas, encuestas\ldots{} si muchas variables se mueven juntas.};
  \end{tikzpicture}
  \vspace{0.1cm}

  \small En todos ellos el PCA responde a la misma pregunta: \alert{¿cuáles son las principales fuentes de variación de mis datos?}
\end{frame}

\begin{frame}{Leer un PCA: primero colorear, después interpretar}
  \begin{columns}[T]
    \begin{column}{0.52\linewidth}
      \begin{block}{Cómo averiguar qué significa una PC}
        \small
        \begin{enumerate}\setlength\itemsep{3pt}
          \item \textbf{Colorea} las muestras según cada metadato.
          \item Calcula la \textbf{R$^2$} de cada PC con cada variable.
          \item Lee las \textbf{cargas}: ¿qué genes definen el eje?
        \end{enumerate}
      \end{block}
    \end{column}
    \begin{column}{0.44\linewidth}
      \begin{block}{Qué puede salir mal}
        \footnotesize
        \begin{itemize}\setlength\itemsep{2pt}
          \item El PCA encuentra la variación \alert{mayor}, no la interesante.
          \item Las variables técnicas (profundidad, lote) pueden dominar: normaliza antes.
          \item Una PC puede mezclar varias cosas (un eje de citoquina y un eje de sexo en una misma PC).
          \item El signo de una PC es arbitrario. Las PCs no son clusters.
          \item Es exploratorio: no se contrasta nada.
        \end{itemize}
      \end{block}
    \end{column}
  \end{columns}
\end{frame}

\begin{frame}{Todas las muestras: domina el tipo celular}
  \centering
  \figfit[0.62]{p3_all_ct_donor}
  \begin{itemize}\small\setlength\itemsep{2pt}
    \item \textbf{PC1} (\numpcaone\,\%): monocitos frente a linfocitos. \textbf{PC2} (\numpcatwo\,\%): células NK frente a linfocitos T.
    \item Dentro de cada grupo de tipo celular los \alert{donantes están mezclados}: los efectos del donante aparecen en PCs posteriores.
  \end{itemize}
\end{frame}

\begin{frame}{Todas las muestras: citoquinas y control vs tratado}
  \centering
  \figfit[0.62]{p3_all_cyt_trt}
  \begin{itemize}\small\setlength\itemsep{2pt}
    \item Los controles PBS (negro) están \alert{dentro} de la nube de muestras tratadas. Unas pocas citoquinas (IL-4) empiezan a salir de ella.
    \item Los efectos de las citoquinas son reales pero \alert{pequeños frente a las diferencias entre tipos celulares}. ¿Significa eso que no hacen nada?
  \end{itemize}
\end{frame}

\begin{frame}{¿Qué captura cada PC?}
  \centering
  \figfit[0.52]{p3_r2}
  \begin{itemize}\small\setlength\itemsep{2pt}
    \item PC1, PC2, PC4: \alert{tipo celular}. PC3, PC5, PC6, PC8: \alert{donante}.
    \item La primera PC con señal de citoquina es la PC3 (R$^2$ 0.13): una PC de baja varianza. Las cuatro primeras PCs acumulan el \numpcafirstfour\,\% de la varianza.
    \item La PC2 se correlaciona con el tamaño de librería solo porque las células NK tienen librerías más pequeñas: es el tipo celular disfrazado.
  \end{itemize}
\end{frame}

\begin{frame}{¿Qué genes definen una PC? Las cargas}
  \centering
  \figfit[0.5]{p3_loadings}
  \begin{itemize}\small\setlength\itemsep{2pt}
    \item \textbf{PC1}: genes de monocitos frente a genes de linfocitos. \textbf{PC2}: genes NK (\gene{GNLY}, \gene{NCAM1}) frente a genes de T vírgenes (\gene{CCR7}, \gene{LEF1}).
    \item La \textbf{PC3} mezcla la \alert{respuesta a interferón} (\gene{IFIT1/2/3}, \gene{ISG15}) con \alert{genes del cromosoma Y} (\gene{UTY}, \gene{USP9Y}): un eje de citoquina y un eje de sexo en una misma PC.
  \end{itemize}
\end{frame}

\begin{frame}{Acercamiento: solo monocitos}
  \centering
  \figfit[0.62]{p3_mono_donor_cyt}
  \begin{itemize}\small\setlength\itemsep{2pt}
    \item Los \alert{donantes} forman grupos claros: las diferencias basales entre personas son grandes. El Donor 2 queda aparte (como en el artículo).
    \item Las citoquinas aún no se separan. Por eso la Parte 4 compara cada citoquina con el PBS del \alert{mismo donante}.
  \end{itemize}
\end{frame}

\begin{frame}{En el artículo: los donantes difieren antes de cualquier tratamiento}
  \centering
  \includegraphics[width=0.95\linewidth,height=0.52\textheight,keepaspectratio]{paper_fig2_donors.png}\\[-1pt]
  {\tiny Oesinghaus \emph{et al.}\ 2025, Fig.\ 2a--g (CC BY-NC-ND 4.0)}
  \begin{itemize}\small\setlength\itemsep{2pt}
    \item Un grupo de donantes (D1, D3, D4, D10) tiene una \alert{firma basal alta de interferón} en PBS (paneles d, e): el ``grupo de interferón'' del artículo; los monocitos del Donor 2 también quedan aparte (Fig.\ 2k), como en nuestro PCA.
    \item Estas diferencias basales son la razón por la que comparamos cada citoquina con el PBS del \alert{mismo donante}.
  \end{itemize}
\end{frame}

\begin{frame}{Monocitos: las PCs posteriores muestran las citoquinas}
  \centering
  \figfit[0.64]{p3_mono_pc34}
  \begin{itemize}\small\setlength\itemsep{2pt}
    \item \textbf{IL-4} sale de la nube hacia valores bajos de PC4, cuyo extremo negativo contiene genes diana de IL-4 (\gene{ALOX15}, \gene{WNT5B}, \gene{FN1}); \textbf{GM-CSF} va hacia valores altos de PC3.
    \item Cada citoquina mueve las muestras en su \alert{propia dirección}, mezclada con los efectos del donante.
  \end{itemize}
\end{frame}

\begin{frame}{Tu turno en el notebook (Parte 3)}
  \small
  Las celdas que empiezan por \code{\# CAMBIA ESTO} contienen parámetros que puedes modificar. Cámbialos, vuelve a ejecutar y mira qué pasa:
  \vspace{0.2cm}

  \begin{tabular}{@{}>{\ttfamily}p{0.25\linewidth}>{\raggedright\arraybackslash}p{0.7\linewidth}@{}}
    \toprule
    \textrm{\textbf{Parámetro}} & \textbf{Prueba\ldots}\\\midrule
    N\_HVG & $200$ y $10\,000$ genes: ¿cuán robusta es la imagen?\\[2pt]
    CYTOKINES\_PANEL & seis citoquinas de la familia de los interferones: ¿cuáles se agrupan?\\[2pt]
    PCA\_CELL\_TYPE & otro tipo celular: ¿el efecto de la citoquina es parecido?\\[2pt]
    COLOR\_BY, PCS & un gen (\gene{UTY}, \gene{STAT1}, \gene{IL1B}), el tamaño de librería, otras PCs (PC3 vs PC4)\\
    \bottomrule
  \end{tabular}
\end{frame}

\begin{frame}{La Parte 3 en una diapositiva}
  \begin{itemize}\setlength\itemsep{7pt}
    \item El PCA es una \alert{rotación} de los datos hacia ejes de varianza decreciente, calculada con la \alert{SVD}: \textbf{puntuaciones} = proyección de las muestras, \textbf{cargas} = pesos de los genes.
    \item Es \alert{no supervisado}: colorear por metadatos, la R$^2$ y las cargas dicen qué significa cada PC.
    \item La \alert{normalización} (Parte 2) decide lo que ve el PCA: sin ella, la PC1 sería la profundidad de secuenciación.
    \item La variación más fuerte es el \alert{tipo celular}; después vienen los efectos del donante; los efectos de las citoquinas son reales pero pequeños y están \alert{mezclados con los del donante}.
    \item $\Rightarrow$ siguiente paso: plantear la pregunta gen a gen, emparejando por donante.
  \end{itemize}
\end{frame}

% ════════════════════════════════════════════════════════════════════════════════════════
\section{Parte 4: Expresión diferencial}

\begin{frame}{Expresión diferencial: la pregunta y el flujo de trabajo}
  \centering
  \textit{Para el gen $g$: ¿es la expresión con la citoquina distinta de la expresión con PBS, más de lo esperable por azar?}\\[0.45cm]
  \begin{tikzpicture}[font=\footnotesize, every node/.style={align=center},
      box/.style={draw=ucoblue,fill=ucoblue!8,rounded corners,minimum height=1.1cm,minimum width=2.6cm,inner sep=3pt}]
    \node[box] (a) {diseño pareado\\ mismo donante,\\ PBS vs citoquina};
    \node[box,right=0.4cm of a] (b) {\textbf{DESeq2}\\ conteos crudos\\ $\Rightarrow$ valor $p$};
    \node[box,right=0.4cm of b] (c) {pruebas múltiples\\ Benjamini--\\ Hochberg};
    \node[box,right=0.4cm of c,fill=accentred!10,draw=accentred] (d) {umbrales\\ $\Rightarrow$ DEG};
    \draw[-{Stealth},thick] (a) -- (b); \draw[-{Stealth},thick] (b) -- (c); \draw[-{Stealth},thick] (c) -- (d);
    \node[font=\scriptsize,text=greyink,below=0.12cm of a] {4.1};
    \node[font=\scriptsize,text=greyink,below=0.12cm of b] {4.2};
    \node[font=\scriptsize,text=greyink,below=0.12cm of c] {4.3};
    \node[font=\scriptsize,text=greyink,below=0.12cm of d] {4.4};
  \end{tikzpicture}
  \vspace{0.4cm}

  \begin{block}{Predice}
    En monocitos CD14 tratados con \IFNg: ¿más genes que \textbf{suben} o que \textbf{bajan}? ¿Decenas, cientos, miles?
  \end{block}
\end{frame}

\begin{frame}{4.1 Diseño pareado: cada donante frente a sí mismo}
  Cada donante se midió con PBS \emph{y} con cada citoquina; $y=\log_2(\text{CPM}+1)$:
  \[ d_{gi}=y^{\text{citoquina}}_{gi}-y^{\text{PBS}}_{gi}\qquad(g=\text{gen},\ i=1,\dots,12\ \text{donantes}),\qquad \lfc_g=\bar d_g=\tfrac1{12}\textstyle\sum_i d_{gi} \]
  \begin{columns}[c]
    \begin{column}{0.55\linewidth}
      \figfit[0.52]{p4_paired}
    \end{column}
    \begin{column}{0.43\linewidth}\small
      \begin{itemize}\setlength\itemsep{3pt}
        \item Cada línea = un donante (rojo: mayor con la citoquina).
        \item \gene{CXCL9}: todos los donantes suben, partiendo de líneas base distintas. $\bar d=\numlfcobs$ ($\times\numfoldobs$).
        \item \gene{ACTB}: las líneas van en ambos sentidos, sin cambio consistente.
      \end{itemize}
    \end{column}
  \end{columns}
\end{frame}

\begin{frame}{4.1 La prueba $t$ pareada, a mano}
  \begin{columns}[c]
    \begin{column}{0.48\linewidth}
      \[
        s_g=\sqrt{\tfrac1{n-1}\textstyle\sum_i(d_{gi}-\bar d_g)^2},\qquad
        t_g=\frac{\bar d_g}{s_g/\sqrt n}
      \]
      \[ p_g=2\,\Pr\bigl(T_{n-1}>|t_g|\bigr),\qquad n=12 \]
      \small
      \gene{CXCL9}, monocitos, \IFNg:
      \begin{itemize}\setlength\itemsep{1pt}
        \item $\bar d=\numlfcobs$, \quad $t=\numtobs$, \quad gl $=11$
        \item $p=\numpobs$
      \end{itemize}
      Bajo H$_0$ (sin efecto) $t_g\sim t_{n-1}$. El valor $p$ es la probabilidad de un $|t|$ al menos así de grande \emph{por azar}.
    \end{column}
    \begin{column}{0.5\linewidth}
      \figfit[0.6]{p4_tdist}
    \end{column}
  \end{columns}
\end{frame}

\begin{frame}{4.2 DESeq2: ¿cómo funciona?}
  \centering
  \begin{tikzpicture}[font=\scriptsize, every node/.style={align=center},
      stp/.style={draw=ucoblue,fill=ucoblue!8,rounded corners,minimum height=2.0cm,minimum width=2.95cm,inner sep=3pt,text width=2.75cm}]
    \node[stp] (a) {\textbf{1. Factores de tamaño}\\[2pt] conteos crudos, \emph{mediana de cocientes}: un CPM robusto que supone que la mayoría de los genes no cambian};
    \node[stp,right=0.35cm of a] (b) {\textbf{2. Modelo}\\[2pt] binomial negativa\\ por gen:\\ \code{\textasciitilde\ donor + cytokine}};
    \node[stp,right=0.35cm of b] (c) {\textbf{3. Dispersión}\\[2pt] ruido de cada gen, \emph{contraído} hacia una tendencia común a todos los genes};
    \node[stp,right=0.35cm of c,fill=accentred!10,draw=accentred] (d) {\textbf{4. Prueba de Wald}\\[2pt] $\log_2$FC del coeficiente de la citoquina, su error estándar, valor $p$};
    \foreach \x/\y in {a/b,b/c,c/d} \draw[-{Stealth},thick] (\x) -- (\y);
  \end{tikzpicture}
  \vspace{0.25cm}

  \begin{columns}[T]
    \begin{column}{0.5\linewidth}\small
      \[ \log\mu_{gi}=\log s_i+\beta_{g,\text{donor}(i)}+\beta_{g,\text{cytokine}(i)} \]
      \begin{itemize}\setlength\itemsep{2pt}
        \item $s_i$: factor de tamaño de la muestra $i$. Se modelan los conteos, no el log-CPM.
        \item \code{donor} da a cada donante su propia línea base: el \alert{diseño pareado}.
      \end{itemize}
    \end{column}
    \begin{column}{0.46\linewidth}\small
      \begin{itemize}\setlength\itemsep{2pt}
        \item Con solo 12 donantes, el ruido propio de un gen se estima mal: compartir información entre genes hace la prueba \alert{más estable}.
        \item En Python: \code{pydeseq2} (una adaptación del paquete de R DESeq2, el estándar del campo).
      \end{itemize}
    \end{column}
  \end{columns}
\end{frame}

\begin{frame}{4.2 DESeq2 por dentro}
  \centering
  \figfit[0.5]{p4_deseq_diag}
  \begin{itemize}\small\setlength\itemsep{2pt}
    \item Los \textbf{factores de tamaño} (izquierda) siguen al tamaño de librería, pero se estiman a partir de los genes que no cambian, no del total.
    \item \textbf{Dispersión} (derecha): el ruido de cada gen (gris) se \alert{acerca a la tendencia} (rojo) para dar el valor que se usa en la prueba (azul).
  \end{itemize}
\end{frame}

\begin{frame}{4.2 DESeq2 frente a nuestra prueba $t$ sobre log-CPM}
  \centering
  \figfit[0.5]{p4_vs_ttest}
  \begin{itemize}\small\setlength\itemsep{2pt}
    \item Tamaños de efecto parecidos ($r=\numreffect$). DEG: DESeq2 \numndegdeseq, prueba $t$ \numndegt, \alert{ambos \numndegboth}.
    \item El $\lfc$ mediano es \nummediandeseq\ con DESeq2 pero \nummediancpm\ con CPM: el \alert{efecto de composición} de la Parte 2, corregido por la mediana de cocientes.
  \end{itemize}
\end{frame}

\begin{frame}{4.3 \numntested\ pruebas: no podemos fiarnos de los valores $p$ crudos}
  \centering
  \figfit[0.42]{p4_pvalues}
  \begin{itemize}\small\setlength\itemsep{2pt}
    \item \numntested\ genes analizados al 5\,\% $\Rightarrow$ $\approx\numnexpectedfp$ falsos positivos \emph{si no pasa nada}.
    \item Experimento nulo (efectos eliminados, signos de los donantes invertidos al azar, prueba $t$): \numnnullraw\ genes con $p<0.05$ y \alert{\numnnulladj} tras la corrección.
    \item Datos reales: \numnrawsig\ genes con $p<0.05$: un pico cerca de 0 sobre un suelo plano.
  \end{itemize}
\end{frame}

\begin{frame}{Benjamini--Hochberg: controlar la tasa de falsos descubrimientos}
  \begin{columns}[T]
    \begin{column}{0.52\linewidth}
      Con $m$ pruebas y $p_{(1)}\le\dots\le p_{(m)}$:
      \[ p^{\mathrm{adj}}_{(i)}=\min_{k\ge i}\Bigl(\tfrac{m}{k}\,p_{(k)}\Bigr)\quad(\le1) \]
      \small
      \alert{FDR} $\le\alpha$: entre los genes declarados significativos con $\mathrm{padj}<0.05$, se espera que como máximo el 5\,\% sean falsos.\\[0.15cm]
      \scriptsize
      \begin{tabular}{lrrrr}\toprule
$i$ & $p_{(i)}$ & $p_{(i)}\,m/i$ & $p^{\mathrm{adj}}_{(i)}$ & sig.\\\midrule
1 & 0.001 & 0.010 & 0.010 & \checkmark\\
2 & 0.008 & 0.040 & 0.040 & \checkmark\\
3 & 0.012 & 0.040 & 0.040 & \checkmark\\
4 & 0.030 & 0.075 & 0.075 & \\
5 & 0.040 & 0.080 & 0.075 & \\
6 & 0.045 & 0.075 & 0.075 & \\
7 & 0.200 & 0.286 & 0.286 & \\
8 & 0.400 & 0.500 & 0.500 & \\
9 & 0.700 & 0.778 & 0.778 & \\
10 & 0.900 & 0.900 & 0.900 & \\
\bottomrule\end{tabular}

    \end{column}
    \begin{column}{0.44\linewidth}
      \figfit[0.62]{p4_bh}\\[-0.1cm]
      \scriptsize monocitos, \IFNg: valores $p$ ordenados y la recta $\alpha k/m$. Todo hasta $k^*=\numkstar$ es significativo.
    \end{column}
  \end{columns}
\end{frame}

\begin{frame}{4.4 ¿Qué llamamos ``diferencialmente expresado''?}
  \begin{columns}[c]
    \begin{column}{0.56\linewidth}
      \figfit[0.78]{p4_volcano}
    \end{column}
    \begin{column}{0.42\linewidth}\small
      \[
        \text{DEG}\;\Longleftrightarrow\;\mathrm{padj}<0.05\;\text{y}\;|\lfc|>0.5
      \]
      \begin{itemize}\setlength\itemsep{3pt}
        \item \textbf{Volcán}: tamaño del efecto ($x$) frente a evidencia ($y$). Significación $\neq$ importancia: pide ambas.
        \item Monocitos + \IFNg: \alert{\numnup\ suben} y \alert{\numndown\ bajan} entre \numntested\ genes analizados.
        \item Etiquetas: la mejor combinación de efecto y evidencia. Los genes clásicos de IFN-$\gamma$ (\gene{CXCL9}, \gene{CXCL10}, \gene{IDO1}, \gene{GBP1}) también son DEG.
      \end{itemize}
    \end{column}
  \end{columns}
\end{frame}

\begin{frame}{4.5 Comparar dos citoquinas}
  \centering
  \figfit[0.48]{p4_compare}\\[-2pt]
  \scriptsize
  \begin{tabular}{@{}lrrrrrr@{}}\toprule
tipo celular & IFN-gamma & IL-4 & ambas & solo IFN-gamma & solo IL-4 & IFN-gamma vs IL-4\\\midrule
monocito CD14 & 2\,609 & 3\,748 & 1\,484 & 1\,125 & 2\,264 & 4\,755\\
linfocito T CD4 & 1\,024 & 1\,617 & 434 & 590 & 1\,183 & 2\,362\\
\bottomrule\end{tabular}

  \begin{itemize}\footnotesize\setlength\itemsep{1pt}
    \item Izquierda: log$_2$FC de cada citoquina vs PBS. Derecha: la comparación directa. La mayoría de los DEG son \alert{específicos} de una citoquina.
  \end{itemize}
\end{frame}

\begin{frame}{En el artículo: las respuestas coinciden con la biología conocida}
  \centering
  \includegraphics[width=0.7\linewidth,height=0.45\textheight,keepaspectratio]{paper_fig1de.png}\\[-1pt]
  {\tiny Oesinghaus \emph{et al.}\ 2025, Fig.\ 1d--e (CC BY-NC-ND 4.0)}
  \begin{itemize}\small\setlength\itemsep{2pt}
    \item Los receptores de citoquinas y las respuestas son específicos de cada tipo celular (d); los genes inducidos (e) son los esperados: IL-4 aumenta \gene{IL4R} e \gene{IL4I1} en linfocitos B, IL-12 aumenta genes inflamatorios y de citotoxicidad en células NK.
    \item El artículo usó \alert{edgeR} sobre pseudobulks con el diseño \code{\textasciitilde\ donor + cytokine}: la misma idea que nuestro análisis con \code{pydeseq2}.
  \end{itemize}
\end{frame}

\begin{frame}{Tu turno en el notebook (Parte 4)}
  \small
  Las celdas que empiezan por \code{\# CAMBIA ESTO} contienen parámetros que puedes modificar. Cámbialos, vuelve a ejecutar y mira qué pasa:
  \vspace{0.2cm}

  \begin{tabular}{@{}>{\ttfamily}p{0.27\linewidth}>{\raggedright\arraybackslash}p{0.68\linewidth}@{}}
    \toprule
    \textrm{\textbf{Parámetro}} & \textbf{Prueba\ldots}\\\midrule
    CYTOKINE, CELL\_TYPE & otra citoquina (IL-4, TNF-alpha, IL-15\ldots) y otro tipo celular: ¿cómo cambia el número de DEG?\\[2pt]
    EXAMPLE\_GENE & un gen que esperes que responda (\gene{ALOX15} para IL-4): ¿cuál obtiene el menor valor $p$?\\[2pt]
    PADJ\_THRESH, LFC\_THRESH & $0.01$ y $1$: ¿cuántos DEG sobreviven?\\[2pt]
    CYTOKINE\_A, CYTOKINE\_B & dos interferones, o dos miembros de la misma familia: ¿qué pareja es la más parecida?\\[2pt]
    CELL\_TYPES\_COMPARE & uno o dos tipos celulares: ¿dónde difieren más las dos citoquinas?\\
    \bottomrule
  \end{tabular}
\end{frame}

\begin{frame}{Salvedades: un modelo didáctico}
  \small
  \begin{tabular}{@{}>{\raggedright\arraybackslash}p{0.2\linewidth}>{\raggedright\arraybackslash}p{0.76\linewidth}@{}}
    \toprule
    \textbf{Diseño} & 12 donantes: la unidad de replicación es el donante. \code{\textasciitilde\ donor + cytokine} elimina las diferencias basales comunes a ambas condiciones, pero los donantes con un tono basal alto de interferón (el ``grupo de interferón'' del artículo) responden de forma distinta.\\[3pt]
    \textbf{Herramienta} & DESeq2 es el estándar, pero no el único: \texttt{edgeR} y \texttt{limma-voom} son alternativas. Todas necesitan \alert{conteos crudos} y réplicas.\\[3pt]
    \textbf{Pseudobulk} & Sumar las células convierte al donante en la réplica. Analizar células únicas como si fueran réplicas infla los falsos positivos.\\[3pt]
    \textbf{Filtrado} & Elimina los genes no analizables \emph{antes} de corregir por pruebas múltiples.\\
    \bottomrule
  \end{tabular}
\end{frame}

\begin{frame}{Mensajes para llevarse a casa}
  \small
  \begin{tabular}{@{}>{\raggedright\arraybackslash}p{0.27\linewidth}>{\raggedright\arraybackslash}p{0.69\linewidth}@{}}
    \toprule
    \textbf{Mira primero} & ¿Qué es un conteo? ¿Cómo se construyó el control? ¿Qué profundidad tiene cada muestra? Los distintos tipos celulares tienen tamaños de librería muy distintos.\\[2pt]
    \textbf{Normaliza} & CPM para las muestras, $\log_2$ para los genes, un filtro para el ruido y z-scores solo para los mapas de calor.\\[2pt]
    \textbf{PCA} & Muestra lo que más varía (tipo celular, después donante); colorea por metadatos y comprueba la R$^2$ y las cargas.\\[2pt]
    \textbf{Contrasta} & DESeq2 sobre conteos crudos, pareado por donante, con corrección por FDR y un umbral de tamaño del efecto.\\[2pt]
    \textbf{Compara} & Dos citoquinas: genes compartidos y específicos, y el contraste directo.\\[2pt]
    \textbf{Siguiente paso} & Convertir las listas de DEG en biología con enriquecimiento funcional (rutas, dianas de factores de transcripción).\\
    \bottomrule
  \end{tabular}
\end{frame}


% ════════════════════════════════════════════════════════════════════════════════════════
\section{Cierre}

\begin{frame}{Toda la práctica en una imagen}
  \centering
  \begin{tikzpicture}[font=\small, every node/.style={align=center},
      box/.style={draw=ucoblue,fill=ucoblue!8,rounded corners,minimum height=1.0cm,minimum width=2.3cm},
      mod/.style={font=\scriptsize\bfseries,text=white,fill=ucoblue,rounded corners,inner sep=2pt}]
    \node[box] (h5) {\code{.h5ad}\\ conteos crudos};
    \node[box,right=0.55cm of h5] (an) {AnnData\\ \numnsamples\ $\times$ \numngenesraw};
    \node[box,right=0.55cm of an] (norm) {CPM, $\log_2$,\\ filtro};
    \node[box,right=0.7cm of norm,yshift=0.85cm] (pca) {PCA\\ colorear por metadatos};
    \node[box,right=0.7cm of norm,yshift=-0.85cm,fill=accentred!10,draw=accentred] (de) {DESeq2, FDR\\ volcán, 2 citoquinas};
    \draw[-{Stealth},thick] (h5) -- (an); \draw[-{Stealth},thick] (an) -- (norm);
    \draw[-{Stealth},thick] (norm) -- (pca); \draw[-{Stealth},thick] (norm) -- (de);
    \node[mod,above=0.1cm of h5] {Parte 1}; \node[mod,above=0.1cm of an] {Parte 2}; \node[mod,above=0.1cm of norm] {Parte 2};
    \node[mod,above=0.1cm of pca] {Parte 3}; \node[mod,below=0.1cm of de] {Parte 4};
  \end{tikzpicture}
  \vspace{0.3cm}

  \begin{itemize}\small\setlength\itemsep{3pt}
    \item Mira los datos antes de calcular nada: ¿qué es un conteo, cómo se construyó el control, qué profundidad tiene cada muestra?
    \item Normaliza con criterio: muestras, genes y mapas de calor son tres problemas distintos.
    \item Deja que el PCA muestre lo que varía; contrasta lo que pida la pregunta, pareando por donante y corrigiendo por pruebas múltiples.
  \end{itemize}
\end{frame}

\begin{frame}{Mini cuestionario: 3 preguntas, 1 minuto}
  \small
  Responde en el Google Form: \alert{[añade aquí el enlace o el código QR del formulario]}
  \vspace{0.15cm}

  \begin{enumerate}\setlength\itemsep{6pt}
    \item La muestra A tiene 5$\times$ más conteos de \emph{cada} gen que la muestra B, solo porque se secuenció a más profundidad. Tras normalizar a CPM, ¿qué esperas para un gen que no cambia?\\
      \textbf{a)} 5$\times$ mayor en A \quad \textbf{b)} igual en A y en B \quad \textbf{c)} 5$\times$ menor en A
    \item En el PCA de todas las muestras, ¿qué separa con más fuerza las muestras (PC1)?\\
      \textbf{a)} la citoquina \quad \textbf{b)} el donante \quad \textbf{c)} el tipo celular \quad \textbf{d)} el tamaño de librería
    \item ¿Por qué usamos valores $p$ ajustados (FDR) cuando analizamos $\sim$13\,000 genes?\\
      \textbf{a)} los valores $p$ crudos se calculan mal \quad \textbf{b)} para limitar los falsos positivos que aparecen por azar al hacer muchas pruebas \quad \textbf{c)} para obtener más genes significativos
  \end{enumerate}
\end{frame}

\appendix

\begin{frame}{Solo para el profesorado: respuestas del cuestionario}
  \small
  \begin{enumerate}\setlength\itemsep{6pt}
    \item \textbf{b)} El CPM divide por el tamaño de librería, así que un gen que no cambia tiene el mismo CPM en ambas muestras (Parte 2).
    \item \textbf{c)} El tipo celular: la PC1 separa monocitos de linfocitos y la PC2 NK de linfocitos T; los efectos de las citoquinas son pequeños en comparación (Parte 3).
    \item \textbf{b)} Con un 5\,\% de significación, $\sim$5\,\% de 13\,000 pruebas ($\approx$ 650 genes) resultan significativas por azar si no pasa nada; la corrección de Benjamini--Hochberg controla la tasa de falsos descubrimientos (Parte 4).
  \end{enumerate}
\end{frame}

\begin{frame}{Referencias}
  \scriptsize
  \begin{itemize}\setlength\itemsep{3pt}
    \item Oesinghaus L.\ \emph{et al.} (2025) A single-cell cytokine dictionary of human peripheral blood. \emph{bioRxiv}. \texttt{doi:10.64898/2025.12.12.693897}
    \item Virshup I.\ \emph{et al.} (2024) anndata: Access and store annotated data matrices. \emph{J.\ Open Source Softw.} 9(101):4371.
    \item Heumos L.\ \emph{et al.} (2023) Best practices for single-cell analysis across modalities. \emph{Nat.\ Rev.\ Genet.} 24:550--572.
    \item Tzec-Interi\'an J.A., Gonz\'alez-Padilla D., G\'ongora-Castillo E.B. (2025) Bioinformatics perspectives on transcriptomics: a comprehensive review of bulk and single-cell RNA sequencing analyses. \emph{Quant.\ Biol.} e78. \texttt{doi:10.1002/qub2.78}
    \item Squair J.W.\ \emph{et al.} (2021) Confronting false discoveries in single-cell differential expression. \emph{Nat.\ Commun.} 12:5692.
    \item Benjamini Y., Hochberg Y. (1995) Controlling the false discovery rate. \emph{J.\ R.\ Stat.\ Soc.\ B} 57:289--300.
    \item Robinson M.D., McCarthy D.J., Smyth G.K. (2010) edgeR. \emph{Bioinformatics} 26:139--140.
    \item Muzellec B.\ \emph{et al.} (2023) PyDESeq2: a python package for bulk RNA-seq differential expression analysis. \emph{Bioinformatics} 39:btad547.
    \item Love M.I., Huber W., Anders S. (2014) Moderated estimation of fold change and dispersion for RNA-seq data with DESeq2. \emph{Genome Biol.} 15:550.
    \item Ritchie M.E.\ \emph{et al.} (2015) limma powers differential expression analyses for RNA-sequencing and microarray studies. \emph{Nucleic Acids Res.} 43:e47.
    \item Pedregosa F.\ \emph{et al.} (2011) Scikit-learn: machine learning in Python. \emph{J.\ Mach.\ Learn.\ Res.} 12:2825--2830.
  \end{itemize}
\end{frame}

\end{document}
